\documentclass[../main.tex]{subfiles}
\graphicspath{{\subfix{../image/}}}
\begin{document}
\chapter{子流形}\label{chapter:ch05-子流形}

许多熟悉的流形自然地作为其他流形的子集出现。我们已经看到,光滑流形的开子集本身可视为光滑流形;但除开集之外,还有许多有趣的例子。本章我们探讨光滑子流形,即作为其他光滑流形之子集的光滑流形。你很快会发现,情况比拓扑子空间的相应理论要微妙得多。

我们首先定义最重要的一类光滑子流形,称为嵌入子流形。它们具有从其包含流形继承来的子空间拓扑,并且恰好就是光滑嵌入的像。正如本章将看到的,它们局部地以 Euclid 空间的线性子空间为模型。由于嵌入子流形最常以光滑映射的水平集形式出现,我们将花一些时间分析水平集何时为光滑子流形。例如,我们将看到常秩映射(特别是光滑淹没)的水平集总是嵌入子流形。

接着,我们引进更一般的一类子流形,称为浸入子流形,它们恰好是单射浸入的像。浸入子流形局部看起来像嵌入子流形,但整体上它的拓扑可能不同于子空间拓扑。

在引进这些基本概念之后,我们处理关于子流形的两个关键技术问题:何时可以把光滑映射的定义域或陪域限制到光滑子流形上而仍保持光滑性?如何把光滑子流形的切空间识别为其外围流形切空间的子空间?然后我们说明子流形理论如何推广到带边子流形的情形。

\subfile{Chapter-05/section-01.tex}
\subfile{Chapter-05/section-02.tex}
\subfile{Chapter-05/section-03.tex}
\subfile{Chapter-05/section-04.tex}
\subfile{Chapter-05/section-05.tex}
\subfile{Chapter-05/section-06.tex}
\end{document}